\documentclass[border=4pt]{standalone}
\usepackage{amsmath,amssymb,mathtools,xcolor,varwidth}
\definecolor{equalteal}{HTML}{00897B}
\definecolor{auxpurple}{HTML}{8E24AA}
\begin{document}
\begin{varwidth}{28em}
\large\bfseries
Now suppose $A$ and $B$ to coincide. By the preceding lemma the angle \textcolor{equalteal}{$\angle dAb$}      vanishes --- the chord and tangent directions merge. Therefore the right lines $Ab$, $Ad$ (which      are always finite) and the intermediate \textcolor{auxpurple}{arc $Acb$} squeeze together and      \textcolor{equalteal}{coincide, becoming equal among themselves}. But the original \textcolor{equalteal}{chord $AB$,      tangent $AD$, and arc $ACB$} are always proportional to these finite three; hence they too must      ultimately stand in the \textcolor{equalteal}{ratio of equality}, $\text{arc}:\text{chord}:\text{tangent}\to 1:1:1$.      This is the crown of the section: as $B\to A$, arc, chord, and tangent become interchangeable,      and every later argument about velocities, forces, and curvature rests upon this single equality.      \textit{Q.E.D.}
\end{varwidth}
\end{document}
